Generative Artificial Intelligence (AI), particularly through Large Language Models (LLMs), has recently emerged as a significant topic in both academic research and industry discussions, driven by its potential to enhance software development practices substantially. Traditionally, software construction involves manual coding, repetitive tasks, and significant clerical work. These activities can be notably improved through Generative AI-driven automation. Tools such as GitHub Copilot and Tabnine have demonstrated practical benefits in automating code generation, improving developer productivity, and reducing cognitive load, as supported by empirical studies examining developers' perceptions and the effectiveness of these tools \cite{sergeyuk2025using}.

However, while automated code generation has received considerable attention, software construction involves a broader set of tasks, including code structuring, refactoring, adapting to evolving requirements, and maintaining architectural coherence and consistency. Despite recognizing the potential benefits of Generative AI, there remains limited clarity regarding its broader capabilities, practical accessibility, maturity levels, limitations, and overall applicability to these software construction tasks beyond mere code generation.

This paper addresses this gap by conducting a Rapid Multivocal Review (RMR) \cite{rocha2025metaprotocol}, mapping the landscape of Generative AI tools and methods explicitly related to software construction. The review aims to provide valuable insights for both academic researchers and industry practitioners seeking to improve their software development processes.